
\documentclass{article}
\usepackage{colortbl}
\usepackage{makecell}
\usepackage{multirow}
\usepackage{supertabular}

\begin{document}

\newcounter{utterance}

\centering \large Interaction Transcript for game `dond', experiment `coop', episode 12 with Qwen3.6{-}35B{-}A3B{-}FP8{-}without{-}reasoning.
\vspace{24pt}

{ \footnotesize  \setcounter{utterance}{1}
\setlength{\tabcolsep}{0pt}
\begin{supertabular}{c@{$\;$}|p{.15\linewidth}@{}p{.15\linewidth}p{.15\linewidth}p{.15\linewidth}p{.15\linewidth}p{.15\linewidth}}
    \# & $\;$A & \multicolumn{4}{c}{Game Master} & $\;\:$B\\
    \hline

    \theutterance \stepcounter{utterance}  
    & & \multicolumn{4}{p{0.6\linewidth}}{
        \cellcolor[rgb]{0.9,0.9,0.9}{
            \makecell[{{p{\linewidth}}}]{
                \texttt{\tiny{[P1$\langle$GM]}}
                \texttt{Jūs spēlējat kooperatīvu sarunu spēli, kurā jums jāvienojas ar citu spēlētāju par to, kā sadalīt objektu kopumu starp jums.} \\
\\ 
\texttt{Noteikumi:} \\
\texttt{(a) Jūs un otrs spēlētājs saņemat objektu kopumu. Katrs no jums saņem arī noslēptu vērtības funkciju, kas norāda, cik jūs vērtējat katru objekta veidu.} \\
\texttt{(b) Jūs apmaiņājaties ziņojumus ar citu spēlētāju, lai vienotos par to, kurš saņem kurus objektus. Katrs var nosūtīt ne vairāk kā 5 ziņojumus vai beigt spēli agri, jebkurā laikā izdarot noslēptu priekšlikumu.} \\
\texttt{(c) Jums lūdz iesniedzt noslēptu priekšlikumu, norādot objektus, kurus vēlaties, formatētu kvadrātiekavās šādi: "{[}Priekšlikums: <skaitlis> <objekta nosaukums>, <skaitlis> <objekta nosaukums>, <...>{]}"} \\
\texttt{(d) Ja jūsu priekšlikumi ir papildinoši, tas ir, ir pietiekami objekti, lai apmierinātu abus priekšlikumus, katrs spēlētājs saņem reitingu, pamatojoties uz saņemto objektu vērtību summu. Pretējā gadījumā jūs abi saņemat nulle punktu.} \\
\texttt{(e) Jūsu mērķis ir maksimizēt sava reitinga un otra spēlētāja reitinga summu.} \\
\\ 
\texttt{Sāksim.} \\
\\ 
\texttt{Pieejamo objektu kopums ir:} \\
\\ 
\texttt{2 gredženi, 1 pulkstenis, 2 klintis} \\
\\ 
\texttt{Jūsu noslēptā vērtības funkcija:} \\
\\ 
\texttt{gredzens: 0, pulkstenis: 4, klints: 3} \\
\\ 
\texttt{SVARĪGI: Jūsu ziņojumi, ja vien tie nav noslēpti priekšlikumi, tiek pārraidīti tieši otram spēlētājam, tāpēc neiekļaujiet nekādu atbildi uz noteikumiem vai tekstu, kas paziņo par jūsu ziņojumu. Lai izdarītu noslēptu priekšlikumu, izmantojiet norādīto formātu. Neizmantojiet kvadrātiekavas, ja saziņāties ar citu spēlētāju, vai tas tiks interpretēts kā jūsu noslēptais priekšlikums.} \\
            }
        }
    }
    & \\ \\

    \theutterance \stepcounter{utterance}  
    & \multicolumn{4}{p{0.6\linewidth}}{
        \cellcolor[rgb]{0.9,0.9,0.9}{
            \makecell[{{p{\linewidth}}}]{
                \texttt{\tiny{[P1$\rangle$GM]}}
                \texttt{Sveiki! Es esmu gatavs sadarbībai.} \\
\\ 
\texttt{Lai sasniegtu mūsu kopīgo mērķi – maksimizēt abē spēlētāju kopējo vērtību – man ir svarīgi uzzināt jūsu vērtējumu par šiem objektiem. Tas palīdzēs mums atrast tādu sadalījumu, kurā kopējā iegūtā vērtība ir pēc iespējas lielāka.} \\
\\ 
\texttt{Lai sāktu, lūdzu, pastāstiet man, cik vērtīgs jums ir katrs objekts (gredzens, pulkstenis, klints)?} \\
            }
        }
    }
    & & \\ \\

    \theutterance \stepcounter{utterance}  
    & & & \multicolumn{2}{p{0.3\linewidth}}{
        \cellcolor[rgb]{0.9,0.9,0.9}{
            \makecell[{{p{\linewidth}}}]{
                \texttt{\tiny{[GM$|$GM]}}
                \texttt{continue} \\
            }
        }
    }
    & & \\ \\

    \theutterance \stepcounter{utterance}  
    & & & \multicolumn{2}{p{0.3\linewidth}}{
        \cellcolor[rgb]{0.9,0.9,0.9}{
            \makecell[{{p{\linewidth}}}]{
                \texttt{\tiny{[GM$|$GM]}}
                \texttt{Sveiki! Es esmu gatavs sadarbībai.} \\
\\ 
\texttt{Lai sasniegtu mūsu kopīgo mērķi – maksimizēt abē spēlētāju kopējo vērtību – man ir svarīgi uzzināt jūsu vērtējumu par šiem objektiem. Tas palīdzēs mums atrast tādu sadalījumu, kurā kopējā iegūtā vērtība ir pēc iespējas lielāka.} \\
\\ 
\texttt{Lai sāktu, lūdzu, pastāstiet man, cik vērtīgs jums ir katrs objekts (gredzens, pulkstenis, klints)?} \\
            }
        }
    }
    & & \\ \\

    \theutterance \stepcounter{utterance}  
    & & \multicolumn{4}{p{0.6\linewidth}}{
        \cellcolor[rgb]{0.9,0.9,0.9}{
            \makecell[{{p{\linewidth}}}]{
                \texttt{\tiny{[GM$\rangle$P2]}}
                \texttt{Jūs spēlējat kooperatīvu sarunu spēli, kurā jums jāvienojas ar citu spēlētāju par to, kā sadalīt objektu kopumu starp jums.} \\
\\ 
\texttt{Noteikumi:} \\
\texttt{(a) Jūs un otrs spēlētājs saņemat objektu kopumu. Katrs no jums saņem arī noslēptu vērtības funkciju, kas norāda, cik jūs vērtējat katru objekta veidu.} \\
\texttt{(b) Jūs apmaiņājaties ziņojumus ar citu spēlētāju, lai vienotos par to, kurš saņem kurus objektus. Katrs var nosūtīt ne vairāk kā 5 ziņojumus vai beigt spēli agri, jebkurā laikā izdarot noslēptu priekšlikumu.} \\
\texttt{(c) Jums lūdz iesniedzt noslēptu priekšlikumu, norādot objektus, kurus vēlaties, formatētu kvadrātiekavās šādi: "{[}Priekšlikums: <skaitlis> <objekta nosaukums>, <skaitlis> <objekta nosaukums>, <...>{]}"} \\
\texttt{(d) Ja jūsu priekšlikumi ir papildinoši, tas ir, ir pietiekami objekti, lai apmierinātu abus priekšlikumus, katrs spēlētājs saņem reitingu, pamatojoties uz saņemto objektu vērtību summu. Pretējā gadījumā jūs abi saņemat nulle punktu.} \\
\texttt{(e) Jūsu mērķis ir maksimizēt sava reitinga un otra spēlētāja reitinga summu.} \\
\\ 
\texttt{Sāksim.} \\
\\ 
\texttt{Pieejamo objektu kopums ir:} \\
\\ 
\texttt{2 gredženi, 1 pulkstenis, 2 klintis} \\
\\ 
\texttt{Jūsu noslēptā vērtības funkcija:} \\
\\ 
\texttt{gredzens: 3, pulkstenis: 0, klints: 2} \\
\\ 
\texttt{SVARĪGI: Jūsu ziņojumi, ja vien tie nav noslēpti priekšlikumi, tiek pārraidīti tieši otram spēlētājam, tāpēc neiekļaujiet nekādu atbildi uz noteikumiem vai tekstu, kas paziņo par jūsu ziņojumu. Lai izdarītu noslēptu priekšlikumu, izmantojiet norādīto formātu. Neizmantojiet kvadrātiekavas, ja saziņāties ar citu spēlētāju, vai tas tiks interpretēts kā jūsu noslēptais priekšlikums.} \\
\\ 
\texttt{Sveiki! Es esmu gatavs sadarbībai.} \\
\\ 
\texttt{Lai sasniegtu mūsu kopīgo mērķi – maksimizēt abē spēlētāju kopējo vērtību – man ir svarīgi uzzināt jūsu vērtējumu par šiem objektiem. Tas palīdzēs mums atrast tādu sadalījumu, kurā kopējā iegūtā vērtība ir pēc iespējas lielāka.} \\
\\ 
\texttt{Lai sāktu, lūdzu, pastāstiet man, cik vērtīgs jums ir katrs objekts (gredzens, pulkstenis, klints)?} \\
            }
        }
    }
    & \\ \\

    \theutterance \stepcounter{utterance}  
    & & & \multicolumn{4}{p{0.6\linewidth}}{
        \cellcolor[rgb]{0.9,0.9,0.9}{
            \makecell[{{p{\linewidth}}}]{
                \texttt{\tiny{[GM$\langle$P2]}}
                \texttt{Paldies par iniciatīvu! Es vēlos sadarbīties, lai mēs abi iegūtu pēc iespējas vairāk.} \\
\\ 
\texttt{Lai mēs varētu atrast optimālu risinājumu, man ir jāzina jūsu vērtējumi. Lūdzu, atbildiet uz manu jautājumu, norādot, cik vērtīgs jums ir katrs objekts (gredzens, pulkstenis, klints).} \\
\\ 
\texttt{Piemēram, ja jums gredzens ir vērtīgs 5 punkti, pulkstenis 10 punkti, bet klints 0 punkti, tad lūdzu, rakstiet: "gredzens: 5, pulkstenis: 10, klints: 0".} \\
\\ 
\texttt{Gaidu jūsu atbildi!} \\
            }
        }
    }
     \\ \\

    \theutterance \stepcounter{utterance}  
    & & & \multicolumn{2}{p{0.3\linewidth}}{
        \cellcolor[rgb]{0.9,0.9,0.9}{
            \makecell[{{p{\linewidth}}}]{
                \texttt{\tiny{[GM$|$GM]}}
                \texttt{continue} \\
            }
        }
    }
    & & \\ \\

    \theutterance \stepcounter{utterance}  
    & & & \multicolumn{2}{p{0.3\linewidth}}{
        \cellcolor[rgb]{0.9,0.9,0.9}{
            \makecell[{{p{\linewidth}}}]{
                \texttt{\tiny{[GM$|$GM]}}
                \texttt{Paldies par iniciatīvu! Es vēlos sadarbīties, lai mēs abi iegūtu pēc iespējas vairāk.} \\
\\ 
\texttt{Lai mēs varētu atrast optimālu risinājumu, man ir jāzina jūsu vērtējumi. Lūdzu, atbildiet uz manu jautājumu, norādot, cik vērtīgs jums ir katrs objekts (gredzens, pulkstenis, klints).} \\
\\ 
\texttt{Piemēram, ja jums gredzens ir vērtīgs 5 punkti, pulkstenis 10 punkti, bet klints 0 punkti, tad lūdzu, rakstiet: "gredzens: 5, pulkstenis: 10, klints: 0".} \\
\\ 
\texttt{Gaidu jūsu atbildi!} \\
            }
        }
    }
    & & \\ \\

    \theutterance \stepcounter{utterance}  
    & & \multicolumn{4}{p{0.6\linewidth}}{
        \cellcolor[rgb]{0.9,0.9,0.9}{
            \makecell[{{p{\linewidth}}}]{
                \texttt{\tiny{[P1$\langle$GM]}}
                \texttt{Paldies par iniciatīvu! Es vēlos sadarbīties, lai mēs abi iegūtu pēc iespējas vairāk.} \\
\\ 
\texttt{Lai mēs varētu atrast optimālu risinājumu, man ir jāzina jūsu vērtējumi. Lūdzu, atbildiet uz manu jautājumu, norādot, cik vērtīgs jums ir katrs objekts (gredzens, pulkstenis, klints).} \\
\\ 
\texttt{Piemēram, ja jums gredzens ir vērtīgs 5 punkti, pulkstenis 10 punkti, bet klints 0 punkti, tad lūdzu, rakstiet: "gredzens: 5, pulkstenis: 10, klints: 0".} \\
\\ 
\texttt{Gaidu jūsu atbildi!} \\
            }
        }
    }
    & \\ \\

    \theutterance \stepcounter{utterance}  
    & \multicolumn{4}{p{0.6\linewidth}}{
        \cellcolor[rgb]{0.9,0.9,0.9}{
            \makecell[{{p{\linewidth}}}]{
                \texttt{\tiny{[P1$\rangle$GM]}}
                \texttt{gredzens: 0, pulkstenis: 4, klints: 3} \\
            }
        }
    }
    & & \\ \\

    \theutterance \stepcounter{utterance}  
    & & & \multicolumn{2}{p{0.3\linewidth}}{
        \cellcolor[rgb]{0.9,0.9,0.9}{
            \makecell[{{p{\linewidth}}}]{
                \texttt{\tiny{[GM$|$GM]}}
                \texttt{continue} \\
            }
        }
    }
    & & \\ \\

    \theutterance \stepcounter{utterance}  
    & & & \multicolumn{2}{p{0.3\linewidth}}{
        \cellcolor[rgb]{0.9,0.9,0.9}{
            \makecell[{{p{\linewidth}}}]{
                \texttt{\tiny{[GM$|$GM]}}
                \texttt{gredzens: 0, pulkstenis: 4, klints: 3} \\
            }
        }
    }
    & & \\ \\

    \theutterance \stepcounter{utterance}  
    & & \multicolumn{4}{p{0.6\linewidth}}{
        \cellcolor[rgb]{0.9,0.9,0.9}{
            \makecell[{{p{\linewidth}}}]{
                \texttt{\tiny{[GM$\rangle$P2]}}
                \texttt{gredzens: 0, pulkstenis: 4, klints: 3} \\
            }
        }
    }
    & \\ \\

    \theutterance \stepcounter{utterance}  
    & & & \multicolumn{4}{p{0.6\linewidth}}{
        \cellcolor[rgb]{0.9,0.9,0.9}{
            \makecell[{{p{\linewidth}}}]{
                \texttt{\tiny{[GM$\langle$P2]}}
                \texttt{Paldies par informāciju! Tagad es zinu jūsu vērtējumus.} \\
\\ 
\texttt{Apkoposim situāciju:} \\
\texttt{{-} **Mani vērtējumi:** gredzens: 3, pulkstenis: 0, klints: 2} \\
\texttt{{-} **Jūsu vērtējumi:** gredzens: 0, pulkstenis: 4, klints: 3} \\
\\ 
\texttt{**Analīze:**} \\
\texttt{Mums ir šādi resursi:} \\
\texttt{{-} 2 gredzeni} \\
\texttt{{-} 1 pulkstenis} \\
\texttt{{-} 2 klintis} \\
\\ 
\texttt{Lai maksimizētu kopējo vērtību, mums vajadzētu sadalīt objektus tā, lai katrs objekts nonāktu pie spēlētāja, kurš tam piešķir augstāku vērtību.} \\
\\ 
\texttt{1. **Gredzeni (2 gab.):**} \\
\texttt{   {-} Man: 3 punkti katrs} \\
\texttt{   {-} Jums: 0 punkti katrs} \\
\texttt{   {-} *Secinājums:* Labāk, ja es saņem abus gredzenus. Kopējā vērtība: 3 + 3 = 6.} \\
\\ 
\texttt{2. **Pulkstenis (1 gab.):**} \\
\texttt{   {-} Man: 0 punkti} \\
\texttt{   {-} Jums: 4 punkti} \\
\texttt{   {-} *Secinājums:* Labāk, ja jūs saņemat pulksteni. Kopējā vērtība: 4.} \\
\\ 
\texttt{3. **Klintis (2 gab.):**} \\
\texttt{   {-} Man: 2 punkti katrs} \\
\texttt{   {-} Jums: 3 punkti katrs} \\
\texttt{   {-} *Secinājums:* Labāk, ja jūs saņemat abas klintis. Kopējā vērtība: 3 + 3 = 6.} \\
\\ 
\texttt{**Ieteicamais sadalījums:**} \\
\texttt{{-} **Es saņemtu:** 2 gredzenus} \\
\texttt{{-} **Jūs saņemtu:** 1 pulksteni, 2 klintis} \\
\\ 
\texttt{**Kopējā vērtība:**} \\
\texttt{{-} Mana daļa: 2 * 3 = 6} \\
\texttt{{-} Jūsu daļa: 4 + 2 * 3 = 10} \\
\texttt{{-} **Kopā:** 16 punkti} \\
\\ 
\texttt{Vai jūs piekrītat šim sadalījumam? Ja jā, lūdzu, apstipriniet, un es iesniegšu noslēpto priekšlikumu.} \\
            }
        }
    }
     \\ \\

    \theutterance \stepcounter{utterance}  
    & & & \multicolumn{2}{p{0.3\linewidth}}{
        \cellcolor[rgb]{0.9,0.9,0.9}{
            \makecell[{{p{\linewidth}}}]{
                \texttt{\tiny{[GM$|$GM]}}
                \texttt{continue} \\
            }
        }
    }
    & & \\ \\

    \theutterance \stepcounter{utterance}  
    & & & \multicolumn{2}{p{0.3\linewidth}}{
        \cellcolor[rgb]{0.9,0.9,0.9}{
            \makecell[{{p{\linewidth}}}]{
                \texttt{\tiny{[GM$|$GM]}}
                \texttt{Paldies par informāciju! Tagad es zinu jūsu vērtējumus.} \\
\\ 
\texttt{Apkoposim situāciju:} \\
\texttt{{-} **Mani vērtējumi:** gredzens: 3, pulkstenis: 0, klints: 2} \\
\texttt{{-} **Jūsu vērtējumi:** gredzens: 0, pulkstenis: 4, klints: 3} \\
\\ 
\texttt{**Analīze:**} \\
\texttt{Mums ir šādi resursi:} \\
\texttt{{-} 2 gredzeni} \\
\texttt{{-} 1 pulkstenis} \\
\texttt{{-} 2 klintis} \\
\\ 
\texttt{Lai maksimizētu kopējo vērtību, mums vajadzētu sadalīt objektus tā, lai katrs objekts nonāktu pie spēlētāja, kurš tam piešķir augstāku vērtību.} \\
\\ 
\texttt{1. **Gredzeni (2 gab.):**} \\
\texttt{   {-} Man: 3 punkti katrs} \\
\texttt{   {-} Jums: 0 punkti katrs} \\
\texttt{   {-} *Secinājums:* Labāk, ja es saņem abus gredzenus. Kopējā vērtība: 3 + 3 = 6.} \\
\\ 
\texttt{2. **Pulkstenis (1 gab.):**} \\
\texttt{   {-} Man: 0 punkti} \\
\texttt{   {-} Jums: 4 punkti} \\
\texttt{   {-} *Secinājums:* Labāk, ja jūs saņemat pulksteni. Kopējā vērtība: 4.} \\
\\ 
\texttt{3. **Klintis (2 gab.):**} \\
\texttt{   {-} Man: 2 punkti katrs} \\
\texttt{   {-} Jums: 3 punkti katrs} \\
\texttt{   {-} *Secinājums:* Labāk, ja jūs saņemat abas klintis. Kopējā vērtība: 3 + 3 = 6.} \\
\\ 
\texttt{**Ieteicamais sadalījums:**} \\
\texttt{{-} **Es saņemtu:** 2 gredzenus} \\
\texttt{{-} **Jūs saņemtu:** 1 pulksteni, 2 klintis} \\
\\ 
\texttt{**Kopējā vērtība:**} \\
\texttt{{-} Mana daļa: 2 * 3 = 6} \\
\texttt{{-} Jūsu daļa: 4 + 2 * 3 = 10} \\
\texttt{{-} **Kopā:** 16 punkti} \\
\\ 
\texttt{Vai jūs piekrītat šim sadalījumam? Ja jā, lūdzu, apstipriniet, un es iesniegšu noslēpto priekšlikumu.} \\
            }
        }
    }
    & & \\ \\

    \theutterance \stepcounter{utterance}  
    & & \multicolumn{4}{p{0.6\linewidth}}{
        \cellcolor[rgb]{0.9,0.9,0.9}{
            \makecell[{{p{\linewidth}}}]{
                \texttt{\tiny{[P1$\langle$GM]}}
                \texttt{Paldies par informāciju! Tagad es zinu jūsu vērtējumus.} \\
\\ 
\texttt{Apkoposim situāciju:} \\
\texttt{{-} **Mani vērtējumi:** gredzens: 3, pulkstenis: 0, klints: 2} \\
\texttt{{-} **Jūsu vērtējumi:** gredzens: 0, pulkstenis: 4, klints: 3} \\
\\ 
\texttt{**Analīze:**} \\
\texttt{Mums ir šādi resursi:} \\
\texttt{{-} 2 gredzeni} \\
\texttt{{-} 1 pulkstenis} \\
\texttt{{-} 2 klintis} \\
\\ 
\texttt{Lai maksimizētu kopējo vērtību, mums vajadzētu sadalīt objektus tā, lai katrs objekts nonāktu pie spēlētāja, kurš tam piešķir augstāku vērtību.} \\
\\ 
\texttt{1. **Gredzeni (2 gab.):**} \\
\texttt{   {-} Man: 3 punkti katrs} \\
\texttt{   {-} Jums: 0 punkti katrs} \\
\texttt{   {-} *Secinājums:* Labāk, ja es saņem abus gredzenus. Kopējā vērtība: 3 + 3 = 6.} \\
\\ 
\texttt{2. **Pulkstenis (1 gab.):**} \\
\texttt{   {-} Man: 0 punkti} \\
\texttt{   {-} Jums: 4 punkti} \\
\texttt{   {-} *Secinājums:* Labāk, ja jūs saņemat pulksteni. Kopējā vērtība: 4.} \\
\\ 
\texttt{3. **Klintis (2 gab.):**} \\
\texttt{   {-} Man: 2 punkti katrs} \\
\texttt{   {-} Jums: 3 punkti katrs} \\
\texttt{   {-} *Secinājums:* Labāk, ja jūs saņemat abas klintis. Kopējā vērtība: 3 + 3 = 6.} \\
\\ 
\texttt{**Ieteicamais sadalījums:**} \\
\texttt{{-} **Es saņemtu:** 2 gredzenus} \\
\texttt{{-} **Jūs saņemtu:** 1 pulksteni, 2 klintis} \\
\\ 
\texttt{**Kopējā vērtība:**} \\
\texttt{{-} Mana daļa: 2 * 3 = 6} \\
\texttt{{-} Jūsu daļa: 4 + 2 * 3 = 10} \\
\texttt{{-} **Kopā:** 16 punkti} \\
\\ 
\texttt{Vai jūs piekrītat šim sadalījumam? Ja jā, lūdzu, apstipriniet, un es iesniegšu noslēpto priekšlikumu.} \\
            }
        }
    }
    & \\ \\

    \theutterance \stepcounter{utterance}  
    & \multicolumn{4}{p{0.6\linewidth}}{
        \cellcolor[rgb]{0.9,0.9,0.9}{
            \makecell[{{p{\linewidth}}}]{
                \texttt{\tiny{[P1$\rangle$GM]}}
                \texttt{Jā, es pilnībā piekrītu šim sadalījumam. Tas ir loģisks un maksimizē mūsu kopējo ieguvumu.} \\
\\ 
\texttt{Es saņemšu 1 pulksteni un 2 klintis.} \\
\texttt{Jūs saņemsiet 2 gredzenus.} \\
\\ 
\texttt{Tā kā mēs esam vienisprātis, es tagad iesniedzu savu noslēpto priekšlikumu:} \\
\\ 
\texttt{{[}Priekšlikums: 1 pulkstenis, 2 klintis{]}} \\
            }
        }
    }
    & & \\ \\

    \theutterance \stepcounter{utterance}  
    & & & \multicolumn{2}{p{0.3\linewidth}}{
        \cellcolor[rgb]{0.9,0.9,0.9}{
            \makecell[{{p{\linewidth}}}]{
                \texttt{\tiny{[GM$|$GM]}}
                \texttt{proposal} \\
            }
        }
    }
    & & \\ \\

    \theutterance \stepcounter{utterance}  
    & & & \multicolumn{2}{p{0.3\linewidth}}{
        \cellcolor[rgb]{0.9,0.9,0.9}{
            \makecell[{{p{\linewidth}}}]{
                \texttt{\tiny{[GM$|$GM]}}
                [0, 1, 2]
            }
        }
    }
    & & \\ \\

    \theutterance \stepcounter{utterance}  
    & & \multicolumn{4}{p{0.6\linewidth}}{
        \cellcolor[rgb]{0.9,0.9,0.9}{
            \makecell[{{p{\linewidth}}}]{
                \texttt{\tiny{[GM$\rangle$P2]}}
                \texttt{Otrs spēlētājs ir iesniedzis slepenu priekšlikumu. Tagad ir jūsu kārta iesniegt savu priekšlikumu.} \\
            }
        }
    }
    & \\ \\

    \theutterance \stepcounter{utterance}  
    & & & \multicolumn{4}{p{0.6\linewidth}}{
        \cellcolor[rgb]{0.9,0.9,0.9}{
            \makecell[{{p{\linewidth}}}]{
                \texttt{\tiny{[GM$\langle$P2]}}
                \texttt{Pamatojoties uz iepriekšējo analīzi, kurā mēs identificējām, ka optimālākais sadalījums maksimizē kopējo vērtību, es iesniedzu šādu priekšlikumu:} \\
\\ 
\texttt{{[}Priekšlikums: 2 gredženi{]}} \\
            }
        }
    }
     \\ \\

    \theutterance \stepcounter{utterance}  
    & & & \multicolumn{2}{p{0.3\linewidth}}{
        \cellcolor[rgb]{0.9,0.9,0.9}{
            \makecell[{{p{\linewidth}}}]{
                \texttt{\tiny{[GM$|$GM]}}
                \texttt{proposal} \\
            }
        }
    }
    & & \\ \\

    \theutterance \stepcounter{utterance}  
    & & & \multicolumn{2}{p{0.3\linewidth}}{
        \cellcolor[rgb]{0.9,0.9,0.9}{
            \makecell[{{p{\linewidth}}}]{
                \texttt{\tiny{[GM$|$GM]}}
                [2, 0, 0]
            }
        }
    }
    & & \\ \\

    \theutterance \stepcounter{utterance}  
    & & & \multicolumn{2}{p{0.3\linewidth}}{
        \cellcolor[rgb]{0.9,0.9,0.9}{
            \makecell[{{p{\linewidth}}}]{
                \texttt{\tiny{[GM$|$GM]}}
                [[0, 1, 2], [2, 0, 0]]
            }
        }
    }
    & & \\ \\

\end{supertabular}
}

\end{document}
